\documentclass[10pt]{article}
\usepackage[margin=0.82in]{geometry}
\usepackage[T1]{fontenc}
\usepackage{lmodern}
\usepackage{amsmath,amssymb,amsthm}
\usepackage{booktabs}
\usepackage{graphicx}
\usepackage{microtype}
\usepackage[hidelinks]{hyperref}
\newtheorem{theorem}{Theorem}
\newtheorem{proposition}{Proposition}
\newcommand{\Z}{\mathbb Z}
\title{An Exhaustive Named-Coordinate Atlas for a Universal Complement Core}
\author{Anonymous}
\date{}
\begin{document}
\maketitle

\begin{abstract}
The common intersection of all complements to a fixed direct factor in an
exact restricted mark cokernel was previously identified as
$8C\cong\mathbb Z/3\oplus\mathbb Z/18$.  We resolve the full incidence of
the sixteen named presentation classes inside this core.  An explicit
$\mathbb Z/9\oplus\mathbb Z/3\oplus\mathbb Z/2$ coordinate model classifies
all $2^{16}=65536$ named supports.  They reach exactly twenty distinct
subgroups, and these are every subgroup of the core.  We give the complete
ten-type distribution by support size and the generating-support polynomial.
There are exactly twenty-five inclusion-minimal generating supports, all
triples containing the ninth named class.  This is a presentation-dependent
atlas, not a canonical coordinate system.
\end{abstract}

\section{Core, scope, and coordinate model}
Let $M$ be the frozen $16\times16$ restricted mark matrix and write
$[e_j]$ for its named coordinate classes.  The predecessor theorem identifies
the universal complement core as
\[
 Q=8C\cong\Z/3\oplus\Z/18,
 \qquad |Q|=54.
\]
We study the labelled family $q_j=8[e_j]$ for $1\leq j\leq16$.

The scope firewall is
\[
 \texttt{NO\_BAD\_EULER\_OR\_ROOT\_NUMBER}.
\]
No canonical Smith coordinates, full table of marks or Burnside ring,
arithmetic or local data, Euler factor, root number, automorphy, or
Hilbert--Polya operator is claimed.

Exact rational residues modulo $M\Z^{16}$ show that
$(q_1,q_3,q_9)$ gives coordinates on
\[
 Q=\Z/9\oplus\Z/3\oplus\Z/2.
\]
In this basis the sixteen rows are
\begin{equation}\label{eq:coordinates}
\begin{split}
&(1,0,0),(6,0,0),(0,1,0),(3,1,0),(0,0,0),(0,0,0),\\
&(4,2,0),(3,2,0),(0,0,1),(0,0,0),(0,1,0),(3,1,0),\\
&(0,0,0),(0,0,0),(2,1,0),(8,2,0).
\end{split}
\end{equation}
For each row, the checker lifts the asserted relation back to an integral
column of $M$.  Thus Equation~\eqref{eq:coordinates} is tied to the original
presentation, not inferred only from element orders.

\section{The complete subgroup atlas}
For $A\subseteq\{1,\ldots,16\}$ put
\[
 Q_A=\langle q_j:j\in A\rangle.
\]
The producer closes all $65536$ supports in the coordinate model.  A separate
queue enumeration starts from zero, adjoins every element of the 54-element
group, and returns its complete subgroup lattice.

\begin{theorem}[Exact lattice coverage]
The family $\{Q_A:A\subseteq\{1,\ldots,16\}\}$ consists of twenty distinct
subgroups and equals the complete subgroup lattice of $Q$.  The type inventory
is
\[
\begin{array}{c|cccccccccc}
\text{type}&1&C_2&C_3&C_6&C_3^2&C_9&C_{18}&C_3{+}C_6&C_3{+}C_9&C_3{+}C_{18}\\
\hline
\#&1&1&4&4&1&3&3&1&1&1.
\end{array}
\]
\end{theorem}

\begin{proof}
Cached subgroup closure applied to Equation~\eqref{eq:coordinates} produces
twenty actual subsets of $Q$.  The independent all-element enumeration also
produces twenty subsets, and exact set comparison gives equality.  GAP
independently enumerates the same ten isomorphism types with the displayed
multiplicities.
\end{proof}

Table~\ref{tab:supports} refines the result by support size.  Its columns use
the type order in the theorem.  Thus it records every one of the $65536$
supports, while keeping supports that generate the same subgroup distinct.

\begin{table}[htbp]
\centering
\scriptsize
\setlength{\tabcolsep}{3.2pt}
\begin{tabular}{r|rrrrrrrrrr|r}
\toprule
$|A|$&$1$&$C_2$&$C_3$&$C_6$&$C_3^2$&$C_9$&$C_{18}$&$C_3{+}C_6$&$C_3{+}C_9$&$Q$&total\\
\midrule
0&1&0&0&0&0&0&0&0&0&0&1\\
1&5&1&6&0&0&4&0&0&0&0&16\\
2&10&5&32&6&13&25&4&0&25&0&120\\
3&10&10&70&32&85&66&25&13&224&25&560\\
4&5&10&80&70&245&95&66&85&940&224&1820\\
5&1&5&50&80&411&80&95&245&2461&940&4368\\
6&0&1&16&50&446&39&80&411&4504&2461&8008\\
7&0&0&2&16&328&10&39&446&6095&4504&11440\\
8&0&0&0&2&165&1&10&328&6269&6095&12870\\
9&0&0&0&0&55&0&1&165&4950&6269&11440\\
10&0&0&0&0&11&0&0&55&2992&4950&8008\\
11&0&0&0&0&1&0&0&11&1364&2992&4368\\
12&0&0&0&0&0&0&0&1&455&1364&1820\\
13&0&0&0&0&0&0&0&0&105&455&560\\
14&0&0&0&0&0&0&0&0&15&105&120\\
15&0&0&0&0&0&0&0&0&1&15&16\\
16&0&0&0&0&0&0&0&0&0&1&1\\
\bottomrule
\end{tabular}
\caption{Complete named-support distribution.  Each row sums to
$\binom{16}{|A|}$.  The penultimate type is $C_3\oplus C_9$ and the last
type column $Q$ is $C_3\oplus C_{18}$.}
\label{tab:supports}
\end{table}

\section{The generating-support complex}
Reading the $Q$ column of Table~\ref{tab:supports} gives the polynomial
\begin{equation}\label{eq:poly}
\begin{split}
P_Q(t)={}&25t^3+224t^4+940t^5+2461t^6+4504t^7+6095t^8\\
&+6269t^9+4950t^{10}+2992t^{11}+1364t^{12}\\
&+455t^{13}+105t^{14}+15t^{15}+t^{16}.
\end{split}
\end{equation}

\begin{proposition}[Minimal named supports]
Exactly twenty-five named supports are inclusion-minimal among those that
generate $Q$.  Every one has size three and contains $S_9$.
\end{proposition}

\begin{proof}
The $\Z/2$ coordinate in Equation~\eqref{eq:coordinates} is nonzero only in
row nine, so every generating support contains $S_9$.  After choosing it,
generation is equivalent to spanning the Frattini quotient of
$\Z/9\oplus\Z/3$ over $\mathbb F_3$.  Exactly twenty-five unordered pairs
of the remaining named rows have nonzero determinant.  Removing any member
then loses either the dyadic coordinate or the required rank, proving
inclusion-minimality.  Exhaustive deletion independently returns the same
twenty-five triples.
\end{proof}

The number three is a property of the named list: no one or two of the
$q_j$ generate $Q$.  It is not the abstract generator rank.  Indeed
$Q\cong\Z/3\oplus\Z/18$ is abstractly two-generated.  Likewise, the
coordinate basis is a verified choice, not a canonical Smith basis.

\section{Reproducibility and conclusion}
The canonical evidence binds the exact C64 matrix and C71 evidence and
manifest.  Producer and checker independently classify all supports and all
abstract subgroups; the checker also verifies every coordinate relation in
the original integer presentation.  GAP reproduces the complete subgroup
type inventory.  Clean replay passes and twenty-eight hostile semantic
mutations are rejected.  C72 thus upgrades the predecessor's list of
twenty-five triples to a complete labelled incidence atlas: every named
support is classified, and every subgroup of the universal core is reached.

\end{document}
